\section{System Design and Methodology}

\subsection{Overall RAG Pipeline}
The system is designed as a citation-grounded Retrieval-Augmented Generation pipeline for Vietnamese labor-law question answering. Given a user question, the assistant retrieves relevant legal evidence from a scoped corpus, assembles the selected evidence, generates an answer, and checks whether the cited legal bases are supported by the retrieved context.

The pipeline is divided into two main stages. The first stage is the offline legal knowledge build. In this stage, selected legal documents are cleaned, structured, enriched with metadata, and indexed. The output of this stage is stored in two complementary forms: a Qdrant hybrid index for text retrieval~\cite{qdrantdocs} and a Neo4j legal graph for legal relationships.

The second stage is the online answering process. When a user submits a question, the system retrieves candidate provisions from Qdrant, expands related evidence through Neo4j when needed, merges and filters the results, and prepares a bounded context for answer generation. The generated answer is then validated before being returned to the user.

Figure~\ref{fig:rag-pipeline} summarizes this overall pipeline. It shows how legal documents are processed offline and how user questions are handled at runtime, from retrieval to citation validation.

\thesisfigure{figures/rag-pipeline.png}{RAG pipeline for citation-grounded Vietnamese labor-law question answering.}{fig:rag-pipeline}

\subsection{Scoped Legal Corpus}
The corpus is intentionally limited to a selected set of Vietnamese labor-law documents. This bounded scope makes the system easier to evaluate and also defines when the assistant should refuse a question or mark it as insufficient-context.

The final corpus contains six document groups: the Labor Code 2019 and the labor-related subset of the Civil Procedure Code 2015, Decree 135/2020/ND-CP and Decree 145/2020/ND-CP, and Circular 09/2020/TT-BLDTBXH and Circular 10/2020/TT-BLDTBXH. These documents were selected because they cover common labor-law issues such as employment contract termination, retirement age, minor workers, working time, and labor dispute procedures.

Vietnamese legal documents are often connected to each other through hierarchy, implementation guidance, and legal references. A code may define the general rule, a decree may provide implementation details, and a circular may contain forms, lists, or technical guidance. Some documents may also clarify, overlap with, or replace earlier rules. Therefore, the corpus cannot be treated as a flat collection of unrelated passages.

Table~\ref{tab:corpus-summary-draft} summarizes the indexed legal corpus used in the system, based on the selected legal documents listed above.

\begin{table}[H]
  \centering
  \caption{Scoped Legal Corpus Summary}
  \label{tab:corpus-summary-draft}
  \small
  \begin{tabularx}{\linewidth}{@{}L{0.34\linewidth}Ycc@{}}
    \toprule
    Document group & Main role in the corpus & Rank & Chunks \\
    \midrule
    Labor Code 2019 & Core labor-law rules, including employment contracts, termination, working time, minor workers, and labor relations. & 1 & 697 \\
    Civil Procedure Code 2015, labor-related subset & Labor dispute procedure and litigation-related provisions. & 1 & 159 \\
    Decree 135/2020/ND-CP & Retirement age roadmap and related implementation details. & 2 & 41 \\
    Decree 145/2020/ND-CP & Detailed guidance for implementing the Labor Code in labor conditions and labor relations. & 2 & 520 \\
    Circular 09/2020/TT-BLDTBXH & Guidance related to minor workers and lists of permitted or restricted work. & 3 & 73 \\
    Circular 10/2020/TT-BLDTBXH & Guidance and forms related to labor contract contents. & 3 & 66 \\
    \midrule
    Total & Six indexed document groups & & 1,556 \\
    \bottomrule
  \end{tabularx}
\end{table}

\subsection{Hierarchy-Aware Legal Processing}
The corpus processing pipeline is designed to preserve the legal structure of Vietnamese normative documents instead of splitting the documents into arbitrary fixed-length text windows. Each legal document is first normalized into a consistent textual format. During this step, the system removes irrelevant formatting noise, standardizes document titles, and separates the main content from headings, tables, appendices, and metadata.

After cleaning, each document is parsed according to its legal hierarchy. The main structural units include document, chapter, section, article, clause, point, and appendix. Each unit is then converted into an evidence chunk. A chunk may correspond to one clause, one group of points under the same clause, or one appendix-level unit, depending on the structure of the original legal text. This strategy keeps legal provisions close to their original legal coordinates while avoiding overly large contexts.

Each evidence chunk contains two groups of information. The first group is retrieval content, which includes the legal text used for semantic and lexical search. The second group is metadata, which records the legal source of the chunk, including document title, document type, article number, clause number, point label, appendix identifier, legal rank, topic labels, actor labels, and issue type. These metadata fields are important for filtering, graph construction, citation display, and citation validation.

Table~\ref{tab:evidence-chunk-structure} summarizes the main fields stored in each legal evidence chunk and their roles in the pipeline.

\begin{table}[htbp]
  \centering
  \caption{Structure of a Legal Evidence Chunk}
  \label{tab:evidence-chunk-structure}
  \small
  \begin{tabularx}{\linewidth}{@{}p{0.22\linewidth}p{0.33\linewidth}X@{}}
    \toprule
    \textbf{Component} & \textbf{Example fields} & \textbf{Purpose in the pipeline} \\
    \midrule
    Legal text & Provision text, shortened excerpt & Used as the main content for dense and sparse retrieval. \\
    Legal hierarchy & Document, chapter, section, article, clause, point, appendix & Preserves the legal coordinates of each provision. \\
    Citation metadata & Document title, article number, clause number, point label, appendix identifier & Supports citation display and citation validation after answer generation. \\
    Legal rank & Law, decree, circular, national-level legal effect & Helps distinguish higher-level rules from implementation guidance. \\
    Semantic labels & Topic, actor, issue type & Supports filtering, graph linking, and query-type-aware evidence selection. \\
    Graph links & Parent relation, cross-reference, implementation guidance, appendix relation & Allows Neo4j expansion to recover legally connected evidence. \\
    \bottomrule
  \end{tabularx}
\end{table}

This representation allows later modules to retrieve evidence by both textual relevance and legal coordinates, while keeping each answer traceable to its original legal source.

This design allows the system to retrieve not only textually similar passages, but also legally identifiable evidence. For example, if a user question mentions a specific article, clause, or document name, the retriever can use the metadata to recover the corresponding legal coordinate. If a retrieved provision belongs to a lower-level document, the system can also inspect its relation to higher-level rules through the legal graph.

\subsection{Hybrid Retrieval Pipeline}
The first retrieval stage is implemented with Qdrant. The system uses a hybrid retrieval strategy that combines dense semantic retrieval and sparse lexical retrieval. Dense retrieval is used to capture semantic similarity between the user question and legal provisions, while sparse retrieval is used to preserve exact matching for article numbers, document names, legal terms, and Vietnamese legal expressions.

For each user question, the backend sends the normalized query to the retrieval module. The dense component retrieves provisions that are semantically close to the question, even when the user uses informal wording. The sparse component gives higher priority to exact legal markers such as \enquote{Điều 35}, \enquote{khoản 1}, \enquote{Bộ luật Lao động 2019}, or document identifiers such as \enquote{Nghị định 145/2020/ND-CP}. The two ranked lists are then merged into a single candidate list.

The system also applies reference recovery when the query explicitly contains a legal coordinate. This step is necessary because exact legal references are highly important in legal QA. If a user asks about a specific article or document, the system should not rely only on semantic similarity. It should directly recover the matching provision and include it in the candidate evidence pool when it exists in the indexed corpus.

The output of this stage is an initial set of candidate legal chunks. These candidates are not immediately sent to the generator. They are first passed to graph expansion, filtering, deduplication, and context budgeting so that the final context remains both legally complete and focused.

\subsection{Policy-Controlled Legal Graph Expansion}
The Neo4j legal graph is used after the initial Qdrant retrieval stage. Its purpose is to recover legally connected evidence that may not be ranked highly by text retrieval alone. The graph represents legal units as nodes and legal relationships as edges. The main node types include document, article, clause, point, appendix, topic, actor, and issue type. The main relationship types include hierarchy relationships, cross-reference relationships, implementation guidance relationships, topic relationships, and appendix-related relationships.

Graph expansion is policy-controlled. The system does not add all neighboring nodes of a retrieved provision. Instead, it decides whether expansion is needed based on the query type and the type of relationship. For example, if a question requires implementation details, the system may expand from a Labor Code article to a related decree. If a question requires a table, form, or list, the system may expand toward an appendix or circular. If a question only asks for a direct definition, graph expansion can be limited to avoid adding distracting evidence.

This policy is important because legal relatedness does not always mean answer relevance. A provision may be connected to the same topic but still not be the correct legal basis for a specific question. Therefore, graph expansion is used as a controlled evidence recovery mechanism, not as an unrestricted neighborhood search.

The expanded evidence is merged with the Qdrant results. Duplicate chunks are removed, and each candidate is kept with its source information, retrieval score, graph relation type, and legal metadata. This makes it possible to inspect why a provision was selected and whether it came from direct retrieval or graph-based expansion.

\subsection{Context Assembly and Budgeting}
After retrieval and graph expansion, the system assembles the final legal context for answer generation. This step is necessary because the generator should receive enough evidence to answer the question, but not so much text that irrelevant provisions reduce answer quality.

The context assembly module first deduplicates overlapping chunks. It then reranks the candidates by combining retrieval relevance, legal-coordinate matching, graph-expansion relevance, legal rank, and query-type compatibility. Evidence that directly matches the user question or its legal coordinate is prioritized. Related evidence from graph expansion is included only when it helps complete the legal basis.

The final context is selected under a fixed budget. This budget limits the number of chunks and the amount of legal text passed to the answer generation module. When the question requires several connected provisions, the system may include a parent article, a related clause, or appendix-level evidence. When the question is simple, the system keeps only the most relevant provisions.

The purpose of context budgeting is to balance completeness and precision. A small context may miss necessary legal evidence, while a large context may introduce provisions that are related but not legally valid for the answer. The system therefore treats context assembly as a separate evidence selection step rather than directly passing all retrieved results to the generator.

\subsection{Citation Validation and Insufficient-Context Handling}
The answer generation module is required to answer only from the selected legal contexts. In the deployed system, a configured language-model provider generates a natural-language answer from the retrieved evidence. For reproducible evaluation, a deterministic extractive generation mode is used so that the benchmark result does not depend on sampling behavior or external model availability.

After generation, the citation validator checks whether the cited legal bases in the answer are present in the selected context. The validator detects three main problems: citations that do not appear in the retrieved evidence, article numbers that are unsupported by the context, and answers that refer to legal sources outside the indexed corpus. If any required citation is missing or unsupported, the answer is treated as invalid for evaluation.

The system also includes insufficient-context handling. If the retrieved evidence is not strong enough to support a legal answer, or if the question requires a document outside the indexed corpus, the assistant should not provide a confident legal conclusion. Instead, it returns an insufficient-context response and explains that the indexed corpus does not contain enough legal basis for the question.

This mechanism is central to the safety of the system. The goal is not only to generate fluent legal explanations, but to ensure that each answer remains traceable to retrieved legal evidence. In this thesis, citation validation does not prove that every legal interpretation is perfect, but it provides a practical control layer against unsupported legal citations and out-of-scope answers.
